\documentclass[a4paper,11pt]{article}
\usepackage[utf8]{inputenc}
\usepackage[T1]{fontenc}
\usepackage[french]{babel}
\usepackage[left=1cm,right=1cm,top=.5cm,bottom=0cm]{geometry}
\usepackage{enumitem}
\usepackage{hyperref}
\usepackage{xcolor}
\usepackage{titlesec}
\usepackage{fontawesome5}
\usepackage{lmodern}
\usepackage[normalem]{ulem}

\definecolor{primary}{RGB}{0, 148, 102}
\definecolor{secondary}{RGB}{80, 80, 80}

\hypersetup{colorlinks=true, linkcolor=primary, filecolor=primary, urlcolor=primary}
\titleformat{\section}{\large\bfseries\color{primary}\uppercase}{}{0em}{}[\titlerule]
\titlespacing{\section}{0pt}{8pt}{5pt}

\begin{document}
\begin{center}
    {\Large \textbf{Lo\"ic Aron MBASSI EWOLO}} \\ \vspace{4pt}
    \textbf{\large Alternant Data Scientist -- IA \& Transition \'Energ\'etique} \\
    \textit{\small Stage INRIA Saclay (Avril--Ao\^ut 2026) -- Disponible en alternance d\`es septembre 2026 (12 mois)} \\ \vspace{4pt}
    \small
    \faMapMarker\ Cergy, France (Mobile IDF) \quad \faPhone\ (+33) 7 59 52 33 98 \quad
    {\color{primary}\faEnvelope\ \href{mailto:loicaron.mbassiewolo@ensea.fr}{loicaron.mbassiewolo@ensea.fr}} \\ \vspace{2pt}
    {\color{primary}\faLinkedin\ \href{https://www.linkedin.com/in/loic-mbassi/}{linkedin.com/in/loic-mbassi}} \quad
    {\color{primary}\faGithub\ \href{https://github.com/Nameless0l}{github.com/Nameless0l}} \quad
    {\color{primary}\faGlobe\ \href{https://mbassiloic.tech}{mbassiloic.tech}}
\end{center}
\vspace{-6pt}
\noindent\small{Mod\'elisation pr\'edictive sur donn\'ees IoT et s\'eries temporelles dans le projet Urban Waste (pr\'ediction de remplissage, optimisation de tourn\'ees, 1\textsuperscript{er} prix national) et pipeline NLP de classification \`a l'INRIA Saclay. Formation double dipl\^ome ENSEA/Polytechnique Yaound\'e en math\'ematiques appliqu\'ees et IA. Travaux de recherche au MILA (Montr\'eal) sur la robustesse et la g\'en\'eralisation des mod\`eles ML. ENGIE, acteur majeur de la transition \'energ\'etique, exploite des capteurs et compteurs intelligents dont l'analyse pr\'edictive peut optimiser la distribution et r\'eduire les pertes.}

\section{Formation}
\begin{itemize}[leftmargin=*, label=\tiny$\bullet$, nosep]
    \item \textbf{ENSEA} \hfill \textit{2025 -- 2027} \\
    \small{Double dipl\^ome d'Ing\'enieur -- Data Science, IA \& Apprentissage Automatique.}
    \item \textbf{\'Ecole Polytechnique de Yaound\'e (1\textsuperscript{\`ere} \'ecole d'ing\'enieur CEMAC)} \hfill \textit{2021 -- 2025} \\
    \small{Dipl\^ome d'Ing\'enieur de Conception (Master 2) : \textbf{Math\'ematiques Appliqu\'ees}, Algorithmique, G\'enie Logiciel.}
\end{itemize}

\section{Comp\'etences Techniques}
\begin{itemize}[leftmargin=*, label=\tiny$\bullet$, nosep]
    \item \textbf{Data Science/ML :} Python, scikit-learn, PyTorch, s\'eries temporelles, pr\'ediction, clustering, classification, hyperparam\`etre tuning.
    \item \textbf{NLP/Deep Learning :} BERT, Transformers, extraction information, classification textes, Computer Vision, architectures CNN/RNN.
    \item \textbf{IoT/Donn\'ees Capteurs :} Traitement s\'eries temporelles, d\'etection anomalies, pr\'ediction maintenance, capteurs intelligents, smart grids.
    \item \textbf{MLOps :} Docker, FastAPI, pipelines Python, CI/CD, monitoring mod\`eles, reproductibilit\'e, notebooks Jupyter.
    \item \textbf{Langues :} Fran\c{c}ais natif, Anglais professionnel (INRIA, recherche scientifique).
\end{itemize}

\section{Exp\'eriences \& Projets}
\begin{itemize}[leftmargin=*, label=\tiny$\bullet$, nosep]
    \item \textbf{Stage INRIA Saclay -- \'Equipe TRIBE (Projet PRIMO)} \hfill \textit{Avril-Ao\^ut 2026} \\
    \textit{Palaiseau (IP Paris)} \hfill \textit{Python, NLP, BERT, Transformers, pandas, scikit-learn}
    \vspace{-2pt}
    \begin{itemize}[leftmargin=0.4cm, nosep, label=\tiny$\bullet$]
        \item \small{Pipeline NLP end-to-end d'analyse 50k+ politiques confidentialit\'e mobiles : tokenization BERT, classification RGPD multi-labels.}
        \item \small{Scoring confiance/transparence interpr\'etable ; implication directe pour analyse donn\'ees smart meters ENGIE.}
        \item \small{Orchestration Python (pandas, scikit-learn, transformers) + Docker ; notebooks reproductibles, versioning Git.}
        \item \small{Collaboration IP Paris/LVMT, programme PEPR MOBIDEC ; r\'edaction scientifique EN.}
    \end{itemize}
    \vspace{3pt}

    \item \textbf{Urban Waste Management -- Smart City (1\textsuperscript{er} Prix National)} \hfill \textit{2024} \\
    \textit{ENSEA} \hfill \textit{ML, IoT, Python, pandas, scikit-learn, s\'eries temporelles, optimisation}
    \vspace{-2pt}
    \begin{itemize}[leftmargin=0.4cm, nosep, label=\tiny$\bullet$]
        \item \small{Mod\'elisation pr\'edictive s\'eries temporelles sur capteurs IoT : pr\'ediction remplissage, -20\% co\^uts collecte.}
        \item \small{Pipeline ML complet : exploration donn\'ees, feature engineering, s\'election mod\`ele, validation crois\'ee, optimisation TSP.}
        \item \small{1\textsuperscript{er} prix CITS National Hackathon (vs 75 \'equipes) ; r\'edaction paper recherche Smart City.}
    \end{itemize}
    \vspace{3pt}

    \item \textbf{Assistant de Recherche @ MILA} \hfill \textit{Juin-Sept 2025} \\
    \textit{Remote (Universit\'e Montr\'eal)} \hfill \textit{Python, PyTorch, TensorFlow, analyses statistiques avanc\'ees}
    \vspace{-2pt}
    \begin{itemize}[leftmargin=0.4cm, nosep, label=\tiny$\bullet$]
        \item \small{\'Etude robustesse et g\'en\'eralisation mod\`eles (Grokking) : implication directe pour fiabilit\'e pr\'edictions consommation.}
        \item \small{Pipelines tests statistiques, visualisations, analyses convergence ; math\'ematiques profondes (alg\`ebre, probas).}
    \end{itemize}
    \vspace{3pt}

    \item \textbf{Syst\`eme Multi-Agents Agricole -- Recommandation Multi-Crit\`eres} \hfill \textit{2024-2025} \\
    \textit{ENSEA} \hfill \textit{Google ADK, RAG, LangChain, Gemini, FastAPI, Docker}
    \vspace{-2pt}
    \begin{itemize}[leftmargin=0.4cm, nosep, label=\tiny$\bullet$]
        \item \small{4 agents IA orchestr\'es avec RAG : fusion donn\'ees h\'et\'erog\`enes (m\'et\'eo, marchés) pour recos multi-crit\`eres.}
        \item \small{Ranking/optimisation multi-objectif : transferable \`a optimisation distribution \'energétique ENGIE.}
        \item \small{D\'eploiement FastAPI + Docker ; gestion conflits, orchestration asynchrone.}
    \end{itemize}
    \vspace{3pt}

    \item \textbf{UROP 2024 -- D\'etection Anomalies ECG (Computer Vision)} \hfill \textit{2024} \\
    \textit{ENSEA} \hfill \textit{TensorFlow, Keras, CNN, s\'eries temporelles, classification}
    \vspace{-2pt}
    \begin{itemize}[leftmargin=0.4cm, nosep, label=\tiny$\bullet$]
        \item \small{Mod\`ele CNN d\'etection anomalies sur s\'eries temporelles ECG : transferable \`a d\'etection anomalies smart meters.}
        \item \small{Transfer learning, data augmentation, \'evaluation m\'etriques cliniques ; collaboration Google DeepMind.}
    \end{itemize}
    \vspace{3pt}

    \item \textbf{IndabaX Africa 2025 -- IA Sant\'e (1\textsuperscript{er} Prix)} \hfill \textit{2025} \\
    \textit{Panafricain} \hfill \textit{Python, pandas, scikit-learn, Flask, Streamlit}
    \vspace{-2pt}
    \begin{itemize}[leftmargin=0.4cm, nosep, label=\tiny$\bullet$]
        \item \small{Classification donn\'ees d\'es\'equilibr\'ees, nettoyage/normalisation ; d\'eploiement API Flask + Streamlit.}
        \item \small{1\textsuperscript{er} prix hackathon panafricain ; exp\'erience d\'eploiement production.}
    \end{itemize}
\end{itemize}

\section{Distinctions \& Leadership}
\begin{itemize}[leftmargin=*, label=\tiny$\bullet$, nosep]
    \item \small{\textbf{1\textsuperscript{er} prix IndabaX Africa 2025} (IA Sant\'e, hackathon panafricain).}
    \item \small{\textbf{1\textsuperscript{er} prix CITS National Hackathon} (Smart City, IoT ML) -- vs 75 \'equipes.}
    \item \small{\textbf{Head of Projects Cell \& Lead AI Cell ENSEA} -- +50\% membres actifs, Google DeepMind.}
    \item \small{\textbf{Certifications :} Supervised ML Stanford/DeepLearning.AI (Oct. 2024), Python for Data Science IBM (Jan. 2024).}
\end{itemize}

\end{document}
